% 16-exhaustiveness.tex — Chapter 16: Pattern Matching Exhaustiveness
\chapter{Pattern Matching Exhaustiveness}
\label{chap:16-exhaustiveness}
\index{exhaustiveness checking}
\index{pattern matching!exhaustiveness}

\pnum
The exhaustiveness checker verifies that every \lain{case} expression or
statement covers all possible values of the scrutinee type.  It is implemented
in \srcref{src/sema/exhaustiveness.h} and called from the resolve pass for
every \cname{STMT\_MATCH} and \cname{EXPR\_MATCH} node.

% -------------------------------------------------------------------------
\section{Wildcard arm}
\label{sec:exh-wildcard}
\index{wildcard arm}
% -------------------------------------------------------------------------

\pnum
A match that contains an \lain{else} arm (the wildcard case) is always
exhaustive.  The wildcard arm has \cname{patterns = NULL} in the AST.  The
checker returns immediately upon encountering such an arm.

% -------------------------------------------------------------------------
\section{Enum exhaustiveness}
\label{sec:exh-enum}
\index{enum exhaustiveness}
% -------------------------------------------------------------------------

\pnum
For enum scrutinees, the checker enumerates all variants in the enum
declaration and verifies that each variant name appears as a pattern in at
least one arm.  A variant is considered covered if:
\begin{itemize}
\item An arm has a pattern that is an \cname{EXPR\_IDENTIFIER} whose name
  matches the variant name, or
\item An arm has a pattern that is an \cname{EXPR\_CALL} (constructor pattern)
  whose callee name matches the variant name.
\end{itemize}

\pnum
If any variant is not covered and there is no wildcard arm, the checker emits an
exhaustiveness error listing the uncovered variants and exits.

% -------------------------------------------------------------------------
\section{Boolean exhaustiveness}
\label{sec:exh-bool}
% -------------------------------------------------------------------------

\pnum
For \lain{bool} scrutinees, the checker verifies that both \lain{true} and
\lain{false} appear as patterns.  A match on \lain{bool} with only one branch
and no wildcard is rejected.

% -------------------------------------------------------------------------
\section{Integer and other types}
\label{sec:exh-other}
% -------------------------------------------------------------------------

\pnum
For integer scrutinees, exhaustiveness cannot generally be proven without a
wildcard arm (because the integer range is too large to enumerate).  The checker
requires a wildcard (\lain{else}) arm for any match on an integer type.

\pnum
For struct types, the checker does not attempt exhaustiveness; a struct is a
single value and any arm that matches a struct literal is sufficient.  A
wildcard arm is still required if the match is not intended to cover every
possible struct layout.
